\documentclass[11pt]{article}

% ---------------------------------------------------------------------------
% Benchmark report: cDFT vs. data-driven surrogates (TabPFN) vs. residual-
% correction models for PC-SAFT interfacial properties.
% Self-contained; compile with pdflatex (no external data needed).
% ---------------------------------------------------------------------------

\usepackage[utf8]{inputenc}
\usepackage[T1]{fontenc}
\usepackage{amsmath,amssymb}
\usepackage{booktabs}
\usepackage{graphicx}
\usepackage{geometry}
\usepackage{multirow}
\usepackage{array}
\usepackage[hidelinks]{hyperref}
\geometry{a4paper,margin=1in}

\newcommand{\gmm}{mN\,m$^{-1}$}
\newcommand{\R}{$R^2$}
\newcommand{\speed}[1]{$\sim\!#1\times$}

\title{\bfseries Physics engine vs.\ machine-learning surrogates for PC-SAFT
interfacial properties:\\ a timing and accuracy benchmark}
\author{Generated benchmark report}
\date{\today}

\begin{document}
\maketitle

\begin{abstract}
We benchmark three families of models that predict interfacial and phase
properties of multicomponent mixtures described by the PC-SAFT Helmholtz energy
functional: (i)~the physics engine itself, classical density functional theory
(cDFT) solved with \texttt{FeOs}; (ii)~\emph{full data-driven} surrogates based
on the TabPFN tabular foundation model (two variants: \texttt{without\_HPO} and
\texttt{with\_HPO}); and (iii)~\emph{residual-correction} models (Gaussian
process regression, GPR; stochastic variational GP, SVGP; and symbolic
regression, SR) that learn the deviation of cDFT from an analytic baseline.
Across all four targets (bubble pressure $P_\mathrm{bubble}$, dew pressure
$P_\mathrm{dew}$, surface tension $\gamma$, and interfacial thickness) the
surrogates reproduce cDFT essentially exactly ($R^2\gtrsim0.9999$ on a held-out
test split) while reducing the per-state-point cost from \emph{seconds} (cDFT)
to \emph{sub-millisecond} (TabPFN, GPR, SR), a speed-up of $10^3$--$10^4\times$.
Accuracy is reported as 5-fold cross-validated (CV) RMSE --- the honest
generalisation estimate, systematically larger than the i.i.d.\ test error. We
place the timings in the context of the GPU-accelerated cDFT of Stierle et~al.\
(2024) and the \texttt{FeOs} reference of Rehner et~al.\ (2023), and show by
direct re-timing on the same node that the multicomponent interface solve ---
\emph{not} the phase-equilibrium flash --- is what makes our per-point cDFT cost
exceed the simple-fluid reference.
\end{abstract}

% ===========================================================================
\section{Introduction and scope}
% ===========================================================================
Classical density functional theory (cDFT) based on the PC-SAFT equation of
state is a quantitatively accurate route to interfacial properties such as
surface tension, but it is computationally expensive: each state point requires
the iterative solution of the Euler--Lagrange equation for the density profile,
dominated by the convolution integrals of the weighted-density functional
\cite{rehner2023feos,stierle2024gpu}. This cost is prohibitive for the
large-scale screening and inverse-design tasks that motivate the field
\cite{stierle2024gpu}.

A machine-learning (ML) surrogate trained on cDFT data can, in principle,
collapse this cost to a single forward pass. We quantify this trade-off for two
surrogate strategies:
\begin{itemize}
  \item \textbf{Full data-driven models} that map the thermodynamic state
        $(T,P,\mathbf{z})$ directly to the target property, using the TabPFN
        tabular foundation model in two variants (\texttt{without\_HPO} and a
        hyperparameter-optimised \texttt{with\_HPO}).
  \item \textbf{Residual-correction models} that predict only the
        \emph{residual} between cDFT and a cheap analytic baseline
        ($\gamma_\mathrm{wsd}$), which the physics already captures most of.
        We use GPR, SVGP, and symbolic regression (SR).
\end{itemize}

\paragraph{Accuracy metric.}
Unless stated otherwise, the headline accuracy figure is the \textbf{5-fold
cross-validated RMSE} (mean $\pm$ standard deviation over folds). CV is reported
because the held-out i.i.d.\ test split is optimistic for this data --- the
folds force the model to generalise across state regions, so the CV RMSE is the
number that matches real deployment error (Sec.~\ref{sec:cv}).

\paragraph{Hardware and environment.}
Three distinct machines are used; cross-machine timings are therefore
order-of-magnitude, not same-rig (Table~\ref{tab:hw}). The cDFT $\gamma$ maps and
the per-stage re-timing (Sec.~\ref{sec:cv}, Table~\ref{tab:stage}) ran with
\texttt{FeOs} (CPU, BLAS-threaded) on the TU~Delft \texttt{c109} \texttt{highmem}
node (2$\times$ AMD~EPYC~9754, 256 cores, $\approx$1.5\,TB RAM). TabPFN
training/inference was timed on an NVIDIA~H100 GPU (Snellius). The residual
models (GPR/SVGP/SR) were timed in the \texttt{py\_A6} environment (PyTorch~2.10
CPU, GPyTorch~1.15.2, scikit-learn~1.6.1) on a local Intel Xeon Silver~4214R @
2.40\,GHz (2 sockets $\times$ 12 cores = 24 threads). All learned models share
the training split $n_\mathrm{train}=13\,552$, $n_\mathrm{test}=2\,904$
(per-property validation split of comparable size); CV is 5-fold on the pooled
data.

\begin{table}[h]
\centering
\caption{Compute hardware used for each part of the benchmark.}
\label{tab:hw}
\small
\begin{tabular}{lll}
\toprule
Machine & Hardware & Used for \\
\midrule
\texttt{c109} (TU~Delft) & 2$\times$ AMD EPYC~9754, 256 cores, $\approx$1.5\,TB RAM (CPU-only) & cDFT maps + per-stage re-timing \\
Snellius                 & NVIDIA H100 GPU & TabPFN training / inference \\
Local workstation        & Intel Xeon Silver~4214R @ 2.40\,GHz, 24 threads & GPR / SVGP / SR residual-model timing \\
\bottomrule
\end{tabular}
\end{table}

% ===========================================================================
\section{Full data-driven surrogates: TabPFN across all targets}
% ===========================================================================

\subsection{Accuracy}
Table~\ref{tab:tabpfn-acc} reports the test/validation accuracy and the 5-fold
CV RMSE of the two TabPFN variants for all four targets. On the i.i.d.\ test
split both variants reproduce cDFT with $R^2\ge0.9998$; \texttt{with\_HPO} is
the more accurate everywhere and is the deployed model. Note the CV RMSE is
markedly larger than the test RMSE (e.g.\ $\gamma$ \texttt{with\_HPO}: CV
$0.040$ vs test $0.006$ \gmm) --- the realistic generalisation cost.

\begin{table}[t]
\centering
\caption{Accuracy of the TabPFN variants on the held-out test and validation
sets, and 5-fold cross-validated RMSE. RMSE units: pressures in bar, $\gamma$ in
\gmm, interfacial thickness in its native unit. RMSE$_\mathrm{CV}$ is the
mean~$\pm$~standard deviation over the 5 folds (same units). Bold marks the
lower CV RMSE in each target group.}
\label{tab:tabpfn-acc}
\small
\begin{tabular}{llrrrrr}
\toprule
Target & Variant & \multicolumn{1}{c}{$R^2_\mathrm{test}$} & \multicolumn{1}{c}{$R^2_\mathrm{val}$}
       & \multicolumn{1}{c}{RMSE$_\mathrm{test}$} & \multicolumn{1}{c}{RMSE$_\mathrm{val}$}
       & \multicolumn{1}{c}{RMSE$_\mathrm{CV}$ (mean\,$\pm$\,std)} \\
\midrule
\multirow{2}{*}{$P_\mathrm{bubble}$}
 & without\_HPO & 0.999930 & 0.999926 & 0.18249 & 0.18249 & $2.732 \pm 0.587$ \\
 & with\_HPO    & 0.999997 & 0.999997 & 0.03582 & 0.03767 & $\mathbf{1.313 \pm 0.302}$ \\
\midrule
\multirow{2}{*}{$P_\mathrm{dew}$}
 & without\_HPO & 0.999947 & 0.999941 & 0.13515 & 0.13812 & $0.298 \pm 0.077$ \\
 & with\_HPO    & 1.000000 & 1.000000 & 0.00214 & 0.00198 & $\mathbf{0.201 \pm 0.184}$ \\
\midrule
\multirow{2}{*}{$\gamma$}
 & without\_HPO & 0.999956 & 0.999961 & 0.04062 & 0.03828 & $0.1630 \pm 0.0172$ \\
 & with\_HPO    & 0.999999 & 0.999999 & 0.00603 & 0.00600 & $\mathbf{0.0395 \pm 0.0094}$ \\
\midrule
\multirow{2}{*}{thickness}
 & without\_HPO & 0.999850 & 0.999823 & 0.00843 & 0.00896 & $\mathbf{0.0272 \pm 0.0143}$ \\
 & with\_HPO    & 0.999896 & 0.999909 & 0.00702 & 0.00642 & $0.0277 \pm 0.0081$ \\
\bottomrule
\end{tabular}
\end{table}

\subsection{Cost: training, inference, memory, model size}
Table~\ref{tab:tabpfn-cost} summarises the micro-benchmark of the two variants
(mean over the four targets). The trade-off is clear:
\begin{itemize}
  \item \texttt{without\_HPO}: cheap to train ($\sim$1.5\,s) \emph{and} to serve
        ($\sim$0.7--1.7\,ms/sample); lightest GPU footprint.
  \item \texttt{with\_HPO}: expensive training ($\sim$17\,min of HPO search) but
        fast, light-weight serving; the deployed $\gamma$ model.
\end{itemize}

\begin{table}[t]
\centering
\caption{TabPFN variant micro-benchmark (mean over the four targets;
$n_\mathrm{train}=13\,552$, $n_\mathrm{test}=2\,904$; H100 GPU).}
\label{tab:tabpfn-cost}
\small
\begin{tabular}{lrr}
\toprule
Metric & without\_HPO & with\_HPO \\
\midrule
Training wall time [s]        & $\sim$1.5      & $\sim$1008 ($\approx$17\,min) \\
Inference [ms/sample]         & 0.7--1.7       & 0.5--5.4 \\
GPU memory [MB]               & 50--86         & $\sim$4900 \\
Model size [MB]               & 87.9           & $\sim$51 \\
\bottomrule
\end{tabular}
\end{table}

% ===========================================================================
\section{cDFT vs.\ the full data-driven surrogate on a $\gamma$ map}
% ===========================================================================
Table~\ref{tab:cdft-vs-ml} compares the cost of generating a $\gamma$ map
($\sim$10\,000 $(T,P,\mathbf{z})$ points per feed) with cDFT against the deployed
\texttt{with\_HPO} TabPFN model predicting the \emph{identical} points. The cDFT
map takes 2--18\,h; the surrogate takes 3.5--6.9\,s, a speed-up of
$10^3$--$10^4\times$ at no meaningful loss of accuracy
($R^2\approx0.9999$, RMSE $\approx0.03$--$0.05$ \gmm).

\begin{table}[t]
\centering
\caption{cDFT (PC-SAFT, CPU/BLAS) vs.\ \texttt{with\_HPO} TabPFN (H100, inference
only) on the same per-feed $\gamma$ map. Speed-up grows with feed because cDFT
cost varies strongly while ML inference is nearly constant.}
\label{tab:cdft-vs-ml}
\small
\begin{tabular}{rrrrrrr}
\toprule
Feed & Points & cDFT map [s] & ML map [s] & Speed-up & \R{} & RMSE [\gmm] \\
\midrule
0 & 9\,987  & 7\,481.1 (2.1\,h)  & 6.89 & 1\,086$\times$  & 0.99997 & 0.035 \\
1 & 9\,764  & 19\,253.2 (5.3\,h) & 5.73 & 3\,359$\times$  & 0.99998 & 0.026 \\
2 & 10\,120 & 18\,286.8 (5.1\,h) & 3.56 & 5\,136$\times$  & 0.99995 & 0.046 \\
3 & 9\,910  & 65\,151.9 (18.1\,h)& 4.34 & 15\,026$\times$ & 0.99997 & 0.032 \\
4 & 10\,050 & 63\,255.5 (17.6\,h)& 3.54 & 17\,874$\times$ & 0.99996 & 0.038 \\
\bottomrule
\end{tabular}
\end{table}

% ===========================================================================
\section{Residual-correction models for $\gamma$}
% ===========================================================================
Instead of learning $\gamma$ directly, the residual models learn the deviation
of cDFT from the analytic baseline $\gamma_\mathrm{wsd}$, so the reconstructed
prediction is $\gamma_\mathrm{pred}=\gamma_\mathrm{base}+\hat{r}$. This keeps the
physics in the loop and leaves only a small, smooth residual to be learned.
Table~\ref{tab:residual} reports their accuracy. GPR (5000-point training
subset) and SVGP (2000 inducing points) target the residual directly (RMSE in
\gmm); SR targets a dimensionless $\varepsilon_\mathrm{base}$ and we quote its
reconstructed $\gamma_\mathrm{cDFT}$-scale RMSE. The CV RMSE columns are
sourced from the cross-validated runs of each model, so test, validation and CV
in a row refer to one and the same fitted model.

\begin{table}[t]
\centering
\caption{Residual-correction models for $\gamma$. RMSE$_\mathrm{CV}$ is the 5-fold
cross-validated mean~$\pm$~std. GPR/SVGP RMSE is on the residual target (\gmm); SR
RMSE is on the reconstructed $\gamma_\mathrm{cDFT}$ scale (\gmm). $R^2_\mathrm{CV}$
is on the same target as the RMSE.}
\label{tab:residual}
\small
\begin{tabular}{llrrrr}
\toprule
Model & Target & RMSE$_\mathrm{test}$ & RMSE$_\mathrm{val}$ & RMSE$_\mathrm{CV}$ (mean\,$\pm$\,std) & $R^2_\mathrm{CV}$ \\
\midrule
GPR (RBF+White, 5000)    & residual   & 0.01584 & 0.01536 & $0.0379 \pm 0.0262$ & 0.9938 \\
SVGP (Mat\'ern-5/2 ARD)  & residual   & 0.03684 & 0.02792 & $0.0233 \pm 0.0030$ & 0.9983 \\
SR (closed-form eq.)     & $\gamma_\mathrm{cDFT}$ & 0.14181 & 0.14235 & $0.2475 \pm 0.0110$ & 0.9983 \\
\bottomrule
\end{tabular}
\end{table}

% ===========================================================================
\section{Cross-validated RMSE for $\gamma$}
\label{sec:cv}
% ===========================================================================
Table~\ref{tab:gamma-comb} collects the 5-fold CV RMSE for $\gamma$ across every
learned model on a common \gmm{} scale. Reading it requires one caveat: TabPFN
predicts $\gamma$ \emph{directly}, whereas GPR and SVGP predict only the small
\emph{residual} $\gamma_\mathrm{cDFT}-\gamma_\mathrm{wsd}$, so their RMSE is on a
quantity of much smaller magnitude --- the residual numbers are therefore not
directly comparable to the direct-$\gamma$ TabPFN numbers, even though both are
in \gmm. Among the full data-driven models, \texttt{with\_HPO} TabPFN is the
best at $0.040$ \gmm. The residual SVGP has the smallest absolute CV RMSE
($0.023$ \gmm), as expected for the easier residual target.

Crucially, the CV RMSE of the deployed \texttt{with\_HPO} model ($0.040$ \gmm)
\emph{agrees with} its deployment RMSE when rolled out over the full cDFT
$\gamma$ map ($0.03$--$0.05$ \gmm, Table~\ref{tab:cdft-vs-ml}), whereas the
i.i.d.\ test RMSE ($0.006$ \gmm) was a factor $\sim$6 optimistic. This is exactly
why we headline the CV figure.

\begin{table}[t]
\centering
\caption{5-fold cross-validated RMSE for $\gamma$, in \gmm\ (mean\,$\pm$\,std).
TabPFN values are on the direct-$\gamma$ target; GPR/SVGP on the residual
target; SR on the reconstructed $\gamma_\mathrm{cDFT}$ scale. The scales differ
(see text); bold marks the best full data-driven model.}
\label{tab:gamma-comb}
\small
\begin{tabular}{llr}
\toprule
Family & Model & RMSE$_\mathrm{CV}$ [\gmm] \\
\midrule
\multirow{2}{*}{Full data-driven}
 & TabPFN without\_HPO & $0.1630 \pm 0.0172$ \\
 & TabPFN with\_HPO    & $\mathbf{0.0395 \pm 0.0094}$ \\
\midrule
\multirow{3}{*}{Residual-correction}
 & GPR (5000)          & $0.0379 \pm 0.0262$ \\
 & SVGP                & $0.0233 \pm 0.0030$ \\
 & SR                  & $0.2475 \pm 0.0110$ \\
\bottomrule
\end{tabular}
\end{table}

% ===========================================================================
\section{Per-state-point timing: the unifying comparison}
% ===========================================================================
Table~\ref{tab:perpoint} is the headline comparison: the wall time to evaluate
$\gamma$ at one state point, for the physics engine and every surrogate.
``Isolated'' is a single-sample call (includes call overhead); ``batched'' is the
amortised per-point cost when predicting the whole map at once (the realistic
deployment cost).

\begin{table}[t]
\centering
\caption{Time per state point for $\gamma$ and speed-up vs.\ cDFT (mean cDFT cost
$3.48$\,s/point). cDFT on \texttt{c109} (AMD EPYC~9754); TabPFN on the NVIDIA
H100; residual models on the local Intel Xeon Silver~4214R CPU. Cross-family
numbers are therefore order-of-magnitude, not same-rig.}
\label{tab:perpoint}
\small
\begin{tabular}{llrr}
\toprule
Model & Mode & Time/point & Speed-up vs.\ cDFT \\
\midrule
cDFT (PC-SAFT)        & per point (range) & 0.75--6.57\,s (mean 3.48\,s) & 1$\times$ \\
\midrule
TabPFN without\_HPO   & isolated & 0.73\,ms  & \speed{4\,700} \\
TabPFN with\_HPO      & isolated & 1.74\,ms  & \speed{2\,000} \\
TabPFN with\_HPO      & batched  & 0.48\,ms  & \speed{7\,200} \\
\midrule
GPR (RBF+White)       & isolated & 0.86\,ms     & \speed{4\,000} \\
GPR (RBF+White)       & batched  & 0.0086\,ms   & \speed{406\,000} \\
SVGP (Mat\'ern-5/2)   & isolated & 4.99\,ms     & \speed{700} \\
SVGP (Mat\'ern-5/2)   & batched  & 0.047\,ms    & \speed{74\,000} \\
SR (closed-form eq.)  & batched  & $\sim$8$\times$10$^{-6}$\,ms & \speed{4\times10^8} \\
\bottomrule
\end{tabular}
\end{table}

\subsection{Training, inference, memory and model size}
Table~\ref{tab:resource} compares the three model families on the same resource
axes as a conventional surrogate benchmark: training cost, per-sample inference
cost, model size, and peak memory. The physics engine (PC-SAFT + cDFT) has no
training phase and no stored model --- its ``inference'' is the full per-point
solve --- whereas the learned models pay a one-off training cost and then serve
predictions at $\sim$1--5\,ms/sample (isolated) or far less when batched.

\begin{table}[t]
\centering
\caption{Resource comparison across the three approaches for $\gamma$. The
full data-driven column lists \texttt{without\_HPO}\,/\,\texttt{with\_HPO} TabPFN
(H100); the hybrid column is WSD\,+\,SVGP residual (1500 inducing points). cDFT
has no training phase or stored model (``---''); \texttt{n/r}~=~not separately
recorded.
$^a$\,cDFT per-point solve, default grid $n_\mathrm{grid}=500$ on \texttt{c109},
range over the 2--8-component feeds; inherently one solve per point (no batching).
$^b$\,isolated single-sample on the local Xeon Silver~4214R; batched deployment
$\approx$0.025--0.047\,ms/sample.
$^c$\,peak resident set size (Slurm \texttt{MaxRSS}) of the cDFT job.
SVGP inference was measured isolated on the local Xeon. TabPFN runs on the H100,
so its constraining memory is VRAM ($\sim$4.9\,GB for \texttt{with\_HPO}); cDFT
and SVGP are CPU-only and use no VRAM.
$^e$\,SVGP training timing/RAM as recorded by the instrumented run on
\texttt{c109} (96 cores): wall 880.2\,s ($\approx$14.7\,min), aggregate CPU
time 83\,293\,s across the cores, and 1724\,MB peak resident set during training.
$^d$\,5-fold CV mean (see Table~\ref{tab:gamma-comb} for $\pm$std); cDFT is the
ground-truth reference (RMSE~$\equiv$~0); SVGP RMSE is on the smaller-magnitude
residual target, not directly comparable to the direct-$\gamma$ TabPFN values.}
\label{tab:resource}
\small
\begin{tabular}{lccc}
\toprule
Metric & PC-SAFT + cDFT & Full data-driven (TabPFN) & WSD + Residual (SVGP) \\
\midrule
$\gamma$ RMSE$_\mathrm{CV}$ [\gmm]  & 0 (reference) & 0.163 / 0.040      & 0.023$^d$ \\
\midrule
Training wall time [s]        & ---           & 1.31 / 1217.4      & 880.2 \\
Training CPU time [s]         & ---           & 8.94 / 2790.7      & 83293$^e$ \\
Per-sample inference [ms]     & 184--1791$^a$ & 0.734 / 1.735      & 5.23$^b$ \\
Model size [MB]               & ---           & 87.9 / 50.4        & 9.1 \\
Host RAM [MB]                 & 622--908$^c$  & 27.6 / 3251.9      & 1724$^e$ \\
GPU memory / VRAM [MB]        & --- (CPU)     & 86.1 / 4853.8      & --- (CPU) \\
\bottomrule
\end{tabular}
\end{table}

% ===========================================================================
\section{Comparison with the literature}
% ===========================================================================
Our cDFT cost should be read against two reference implementations.

\paragraph{Where the per-point seconds actually go.}
To make our per-point cost comparable to literature figures (which time a single
simple-fluid interface), we re-ran the V5 contour-map sweep for the five study
feeds on the same \texttt{c109} node, instrumented to separate the
\emph{phase-equilibria} stage (critical point + bubble/dew envelope, once per
feed) from the \emph{cDFT interfacial sweep} (TP-flash + planar-interface solve
at every $(T,P)$ point), at the default grid ($n_\mathrm{grid}=500$, no
critical-region enhancement). Table~\ref{tab:stage} gives the result: phase
equilibria is only $\approx$2.5\% of the stage time; the \emph{cDFT interface
solve dominates ($\approx$97.5\%)}. The per-point cost moreover broadly rises
with the number of components --- from $0.18$\,s/point for the 2-component feed
to $1.8$\,s/point for the 8-component feed, a $\sim$10$\times$ spread driven by
the coupled multicomponent density profiles, not by the flash --- though it is
not strictly monotonic (the 4-component feed at $0.60$\,s/point exceeds the
6-component feed at $0.53$\,s/point), since per-point cost also depends on how
near-critical and hard-to-converge a feed is.
(These absolute times sit below the production-grid maps of
Table~\ref{tab:cdft-vs-ml} because the latter used the finer $n_\mathrm{grid}=2048$
grid; the $\approx$2.5\% phase-equilibria fraction and the component-count
scaling are the grid-independent take-aways.)

\begin{table}[t]
\centering
\caption{Re-timed cDFT $\gamma$ map on \texttt{c109} (2$\times$ AMD EPYC~9754;
default grid $n_\mathrm{grid}=500$, hybrid $4\times4$ threading):
phase-equilibria stage vs.\ cDFT interfacial sweep, per feed. ``Total/pt'' is the
combined stage time per $(T,P)$ point.}
\label{tab:stage}
\small
\begin{tabular}{rrrrrrrr}
\toprule
Feed & Comp.\ & Points & Phase-eq [s] & cDFT sweep [s] & cDFT/pt [s] & Total/pt [s] & Phase-eq share \\
\midrule
0 & 2 & 95  & 1.6 & 15.8  & 0.167 & 0.184 & 9.2\% \\
1 & 4 & 100 & 1.8 & 57.9  & 0.579 & 0.597 & 3.0\% \\
2 & 6 & 100 & 2.2 & 51.1  & 0.511 & 0.533 & 4.1\% \\
3 & 7 & 105 & 3.5 & 181.1 & 1.725 & 1.758 & 1.9\% \\
4 & 8 & 100 & 3.1 & 176.1 & 1.761 & 1.791 & 1.7\% \\
\midrule
\multicolumn{2}{l}{Total / aggregate} & 500 & 12.2 & 481.9 & --- & --- & \textbf{2.5\%} \\
\bottomrule
\end{tabular}
\end{table}

\paragraph{\texttt{FeOs} reference (Rehner et~al.\ 2023).}
The \texttt{FeOs} framework \cite{rehner2023feos} --- the same engine used to
generate our $\gamma$ maps --- reports (1D, 1024 grid points, Intel Core
i7-7700K) a single planar surface tension in $0.09$\,s and a 50-composition
surface-tension diagram in $14.2$\,s ($\approx0.28$\,s/composition), with one
Helmholtz-energy evaluation (with hyper-dual derivatives) at $\sim$0.9--1.5\,$\mu$s.
The gap to our $0.75$--$6.6$\,s/point is therefore \emph{not} explained by the
phase-equilibrium flash (which Table~\ref{tab:stage} shows is negligible,
$\approx$2.5\%): it is that our ``point'' is a coupled \emph{multicomponent}
(2--8 species) interface solve --- intrinsically 10--70$\times$ more expensive
than the simple 1D pure-/binary-fluid interface that the reference times --- on
a finer grid. The order of magnitude is consistent once this is accounted for.

\paragraph{GPU-accelerated 3D cDFT (Stierle et~al.\ 2024).}
Even the state-of-the-art GPU-accelerated cDFT of Stierle et~al.\
\cite{stierle2024gpu} (JAX, backward-mode automatic differentiation, RTX~4090)
requires \emph{seconds} per calculation: one 3D equilibrium density profile takes
$0.42$\,s ($N_\mathrm{grid}=8$) to $21.8$\,s ($N_\mathrm{grid}=256$), and a full
3D adsorption isotherm (20 pressures, $N_\mathrm{grid}=64$) takes $\approx16$\,s.
Their GPU implementation is $\sim$50--60$\times$ faster than a single CPU core and
up to $550\times$ faster than the \texttt{FeOs} CPU reference at
$N_\mathrm{grid}=128$ (where \texttt{FeOs} needs $1580$\,s).
Table~\ref{tab:lit} places our numbers alongside these.

\begin{table}[t]
\centering
\caption{Literature context for the cost of a cDFT evaluation vs.\ our ML
surrogate. Tasks differ (1D interface vs.\ 3D adsorption profile), so this is a
qualitative placement, not a controlled comparison.}
\label{tab:lit}
\small
\begin{tabular}{lll}
\toprule
Source / method & Task & Wall time \\
\midrule
This work, cDFT (\texttt{FeOs}, c109) & 1D $\gamma$/point, 2--8 comp, $n_\mathrm{grid}{=}2048$ & 0.75--6.6\,s \\
This work, cDFT (\texttt{FeOs}, c109) & 1D $\gamma$/point, 2--8 comp, $n_\mathrm{grid}{=}500$ (re-timed) & 0.18--1.8\,s \\
Rehner 2023 \cite{rehner2023feos}, \texttt{FeOs} CPU & 1D surface tension, 1024 pts & 0.09\,s \\
Rehner 2023 \cite{rehner2023feos}, \texttt{FeOs} CPU & 50-composition $\sigma$ diagram & 14.2\,s \\
Stierle 2024 \cite{stierle2024gpu}, JAX GPU & 3D density profile ($N{=}64$)   & 0.51\,s \\
Stierle 2024 \cite{stierle2024gpu}, JAX GPU & 3D density profile ($N{=}256$)  & 21.8\,s \\
Stierle 2024 \cite{stierle2024gpu}, JAX GPU & 3D isotherm (20 $p$, $N{=}64$)  & $\approx$16\,s \\
Stierle 2024 \cite{stierle2024gpu}, \texttt{FeOs} CPU & 3D profile ($N{=}128$) & 1580\,s \\
\textbf{This work, TabPFN/GPR/SR surrogate} & 1D $\gamma$ per state point & \textbf{$5\times10^{-4}$--$1.7\times10^{-3}$\,s} \\
\bottomrule
\end{tabular}
\end{table}

\paragraph{Take-away.}
GPU acceleration brings cDFT from minutes/hours to \emph{seconds} per calculation
\cite{stierle2024gpu}; our ML surrogates bring it a further $10^3$--$10^4\times$
down, to \emph{sub-millisecond} per point --- i.e.\ faster than even
GPU-accelerated cDFT --- while preserving cDFT-level accuracy
($R^2\gtrsim0.9999$). The two approaches are complementary: fast (GPU) cDFT
generates the training data; the surrogate then serves predictions at a cost
compatible with high-throughput screening and gradient-based inverse design.

% ===========================================================================
\section{Conclusions}
% ===========================================================================
\begin{itemize}
  \item Across all four targets, both TabPFN variants reproduce cDFT with
        $R^2\ge0.9998$ on the test split; \texttt{with\_HPO} is the more accurate
        and the right choice for deployment (CV $\gamma$ RMSE $=0.040$ \gmm).
  \item Cross-validation, not the i.i.d.\ test split, is the honest accuracy
        metric: the CV $\gamma$ RMSE ($0.040$ \gmm) matches the deployment error
        on the full map ($0.03$--$0.05$ \gmm), while the test RMSE ($0.006$ \gmm)
        was $\sim$6$\times$ optimistic.
  \item Replacing cDFT by the surrogate accelerates a $\gamma$ map by
        $10^3$--$10^4\times$ (hours $\to$ seconds) at negligible accuracy cost.
  \item Residual-correction models also work well; on the residual target SVGP
        has the smallest CV RMSE ($0.023$ \gmm) and GPR/SR are the cheapest to
        serve ($\sim$8\,$\mu$s and $\sim$8\,ns per point batched).
  \item Direct re-timing on \texttt{c109} shows the cDFT per-point cost is
        dominated by the multicomponent interface solve ($\approx$97.5\%), not the
        phase-equilibrium flash ($\approx$2.5\%); per-point cost scales with
        component count ($0.18\to1.8$\,s for 2$\to$8 species).
  \item Per state point, the learned models cost sub-millisecond vs.\ seconds for
        cDFT --- faster than even the GPU-accelerated cDFT of the literature
        \cite{stierle2024gpu}, while the physics engine (\texttt{FeOs}
        \cite{rehner2023feos}) remains the source of ground truth.
\end{itemize}

% ===========================================================================
\begin{thebibliography}{9}

\bibitem{rehner2023feos}
P.~Rehner, G.~Bauer, and J.~Gross,
\newblock FeOs: An Open-Source Framework for Equations of State and Classical
Density Functional Theory,
\newblock \emph{Ind.\ Eng.\ Chem.\ Res.} \textbf{62} (2023) 5347--5357.
\newblock doi:10.1021/acs.iecr.2c04561.

\bibitem{stierle2024gpu}
R.~Stierle, G.~Bauer, N.~Thiele, B.~Bursik, P.~Rehner, and J.~Gross,
\newblock Classical density functional theory in three dimensions with
GPU-accelerated automatic differentiation: Computational performance analysis
using the example of adsorption in covalent-organic frameworks,
\newblock \emph{Chem.\ Eng.\ Sci.} \textbf{298} (2024) 120380.
\newblock doi:10.1016/j.ces.2024.120380.

\end{thebibliography}

\end{document}
